% !PRINTED 483
\weberpage{483}
If $-i\omega > \sqrt2$, then $f_1^2(\omega) > \sqrt2$ as well; for as
$-i\omega$ increases from $\sqrt2$ to $\infty$, $f_1(\omega)^2$ likewise
increases monotonically\footnote{For from § 54, (11) there follows, by the substitution $(\omega, 1-\frac1{\omega})$:
\[ d\,f_1(\omega)^8 = -\frac{\pi i}{3}\vartheta_{10}^4\left[f(\omega)^8+f_1(\omega)^8\right]d\omega. \]},
from $\sqrt2$ to $\infty$, and consequently $A_1$ is negative for
$r = 0, 2, 4$.
The first factor $A_1+1$ vanishes, as appears from the already
known result (24), for $r = 4$; hence the third factor
$A_1^3+2A_1^2+A_1+4$ vanishes for $r = 0$, while
$A_1+2$ vanishes for $r = 2$.
We thus obtain from (25):
\weq{27}{f_1(\sqrt{-10})^2 = \sqrt2\,y, \quad y - \frac1y = 1,}
\weq{28}{f_1(\sqrt{-22})^2 = \sqrt2\,y, \quad y - \frac1y = 2,}
\weq{29}{f_1(\sqrt{-26})^2 = \sqrt2\,y, \quad y^6-2y^5-2y^4+2y^2-2y-1=0.}

\begin{center}\rule{3cm}{0.4pt}\end{center}

IV. As an example of another kind of use of the
\gesperrt{Schlaefli} modular equations, one which leads to the otherwise hardly
 accessible class invariant $f(\sqrt{-41})$, let the following
serve.
We set:
\weq{30}{2n\omega^2 + 2r\omega + n = 0,}
\weq{31}{\omega = \frac{-r+\sqrt{-m}}{2n}, \quad m = 2n^2-r^2.}
Here $r$ may denote any number which makes $m$ positive,
but which has no divisor in common with $n$ [since otherwise
(30) is imprimitive]. Then, by § 34:
\weq{32}{f_2\left(\frac{\omega}{n}\right) = f_1[2(n\omega+r)] = \frac{e^{-\frac{r\pi i}{12}}\sqrt2}{f_2(n\omega)},}
\[ f_1[2(n\omega+r)] = f_1(r+\sqrt{-m}) = e^{-\frac{r\pi i}{24}}x, \]
%CONT>>
